\documentclass{article}
\usepackage[T1]{fontenc}
\usepackage{amssymb, amsmath, graphicx, subfigure}
\usepackage{hyperref}

\setlength{\oddsidemargin}{.25in}
\setlength{\evensidemargin}{.25in}
\setlength{\textwidth}{6in}
\setlength{\topmargin}{-0.4in}
\setlength{\textheight}{8.5in}

\newcommand{\heading}[6]{
  \renewcommand{\thepage}{#1-\arabic{page}}
  \noindent
  \begin{center}
  \framebox{
    \vbox{
      \hbox to 5.78in { \textbf{#2} \hfill #3 }
      \vspace{4mm}
      \hbox to 5.78in { {\Large \hfill #6  \hfill} }
      \vspace{2mm}
      \hbox to 5.78in { \textit{Instructor: #4 \hfill #5} }
    }
  }
  \end{center}
  \vspace*{4mm}
}

\newtheorem{theorem}{Theorem}
\newtheorem{definition}[theorem]{Definition}
\newtheorem{remark}[theorem]{Remark}
\newtheorem{lemma}[theorem]{Lemma}
\newtheorem{corollary}[theorem]{Corollary}
\newtheorem{proposition}[theorem]{Proposition}
\newtheorem{claim}[theorem]{Claim}
\newtheorem{observation}[theorem]{Observation}
\newtheorem{fact}[theorem]{Fact}
\newtheorem{assumption}[theorem]{Assumption}

\newenvironment{proof}{\noindent{\bf Proof:} \hspace*{1mm}}{
	\hspace*{\fill} $\Box$ }
\newenvironment{proof_of}[1]{\noindent {\bf Proof of #1:}
	\hspace*{1mm}}{\hspace*{\fill} $\Box$ }
\newenvironment{proof_claim}{\begin{quotation} \noindent}{
	\hspace*{\fill} $\diamond$ \end{quotation}}

\newcommand{\problemset}[3]{\heading{#1}{CMSC 27100-3}{#2}{Robert Rand}{#3}{Problem Set #1}}



%%%%%%%%%%%%%%%%%%%%%%%%%%%%%%%%%%%%%%%%%%%%%%%%%%%%%%%%%%%%%%%%%%%%%%%%%%%%%%%
% PLEASE MODIFY THESE FIELDS AS APPROPRIATE
\newcommand{\problemsetnum}{2}          % problem set number
\newcommand{\duedate}{Due: Friday, October 15, 2021}  % problem set deadline
\newcommand{\studentname}{}    % full name of student (i.e., you)
% PUT HERE ANY PACKAGES, MACROS, etc., ADDED BY YOU
% ...
%%%%%%%%%%%%%%%%%%%%%%%%%%%%%%%%%%%%%%%%%%%%%%%%%%%%%%%%%%%%%%%%%%%%%%%%%%%%%%%


%%%%%%%%%%%%%%%%%%%%%%%%%%%%%%%%%%%%%%%%%%%%%%%%%%%%%%%%%%%%%%%%%%%%%%%%%%%%%%%
\begin{document}
\problemset{\problemsetnum}{\duedate}{\studentname}



\section{Predicate Calculus II}

\noindent{\bf{Problem 1:}} Assume $\forall x, P x \wedge Q x$. Prove that $\forall x, Q x \vee P x$.\\[10pt]

\noindent{\bf{Problem 2:}} Assume $\forall x, \forall y, P x \to Q y$ and $\exists z, \neg Q z$. Prove that $\exists x, \neg P x$. \\[10pt]


\section{Functions and relations}

\noindent{\bf{Problem 3}:} Write a bijection $f$ between the sets
$$ \{ (x,y) \in \mathbb{N} \times \mathbb{N} \mid x = y + 1 \vee y = x + 1) \} $$ 
and $\mathbb{Z}$. Also write the inverse function.\\[10pt]

\noindent{\bf{Problem 4}:} Let $R$ be a relation over sets
where $(A,B) \in R$ means ``there is a bijection from $A$ to $B$''. For each of the questions below include a proof (a proof sketch suffices for 3).
\begin{enumerate}
\item Is $R$ reflexive, irreflexive or neither?
\item Is $R$ symmetric, asymmetric, antisymmetric or none?
\item Is $R$ transitive?
\end{enumerate}~\\[10pt]

\noindent{\bf{Problem 5}:} Let $R$ be the \emph{divides} relation restricted to pair of naturals ($R \subseteq \mathbb{N} \times \mathbb{N}$). For each of the questions below include a proof.
\begin{enumerate}
\item Is $R$ reflexive, irreflexive or neither?
\item Is $R$ symmetric, asymmetric, antisymmetric or none?
\item Is $R$ transitive?
\end{enumerate}


\end{document}
%%%%%%%%%%%%%%%%%%%%%%%%%%%%%%%%%%%%%%%%%%%%%%%%%%%%%%%%%%%%%%%%%%%%%%%%%%%%%%%